\begin{claim}{D.2.50}{Proved}{The integral bottom square is nonzero}
Let $L=Y\langle7\rangle$, and retain the specified classes
$a_7\in H^7(L;\ZZ/4)$ and $b_8=\beta_4a_7\in H^8(L;\ZZ)$.
Then
\[
              0\ne b_8^2\in H^{16}(L;\ZZ),
              \qquad 2b_8^2=0.
              \tag{D.2.50a}
\]
There is an integral class $\omega\in H_{15}(L;\ZZ)$ such that
\[
 \big\langle a_7\smile\operatorname{red}_4b_8,
                     \operatorname{red}_4\omega\big\rangle=2,
 \qquad
 \big\langle\theta,\operatorname{red}_2\omega\big\rangle=1.
 \tag{D.2.50b}
\]
Here $\theta$ is the half-product class of D.2.45, with its
stated $\operatorname{Sq}^1$ ambiguity. The class $\omega$ is
torsion of order at least four. The argument does not show that
$\omega$ is in the stable Hurewicz image; consequently it does
not settle the stable-spherical pairing or the actual composite
$h_Ml_4$.
\end{claim}
\begin{proof}
\textbf{An actual integral cobar model.}
Choose a six-reduced simplicial model $K$ of the six-connected
space $L$. Such a model can be obtained from a minimal Kan
model: minimality and the vanishing homotopy groups through
degree six make all simplices in those degrees degenerate.
The same model-existence statement is recorded explicitly in
\cite[proof of Lemma 5.4.2, p.~43 in the opened author version]
{I_ABFJ2022}. Let $C=C_*(K;\ZZ)$ be normalized chains, with the
Alexander--Whitney coproduct. Thus $\overline C_i=0$ for
$i\le6$. The canonical integral cobar construction is
\[
 \Omega C=T(s^{-1}\overline C),\qquad
 d_\Omega=d_1+d_2,
 \quad d_1=-s^{-1}d_Cs,
 \quad d_2=(s^{-1}\otimes s^{-1})\overline\Delta s.
 \tag{D.2.50c}
\]
It is naturally quasi-isomorphic as a chain algebra to the
integral chains of the Kan loop group, hence computes the
actual Pontryagin algebra of $\Omega L$; see
\cite[pp.~3 and 5]{FGHHessParentScottTonks2005}.
The coproduct in \textup{(D.2.50c)} is the ordinary
Alexander--Whitney coproduct. Its dual therefore computes the
ordinary cup product on $L$.

Choose integral cochains $A\in C^7$ and $B\in C^8$ with
\[
       A\bmod4\text{ representing }a_7,
       \qquad dA=4B,
       \qquad dB=0.
       \tag{D.2.50d}
\]
Then $B$ represents $b_8$. Choose a cycle $x\in C_7$
representing a generator of $H_7(L;\ZZ)=\ZZ/4$, and
$y\in C_8$ with $d_Cy=4x$. The integer $A(x)$ is odd, and
\[
                         B(y)=A(x).
                         \tag{D.2.50e}
\]
Both $x$ and $y$ have zero reduced coproduct by degree.
Put $u=s^{-1}x$ and $v=-s^{-1}y$. Then
$d_\Omega u=0$ and $d_\Omega v=4u$. D.2.49 and its
chain-choice independence imply
\[
       2(uv-vu)=d_\Omega W
                    \quad\text{for some }W\in(\Omega C)_{14}.
       \tag{D.2.50f}
\]
In particular this is a boundary in the canonical model,
not only after projecting to another loop space.

\textbf{The nullchain gives an integral fifteen-cycle.}
Every positive generator of $\Omega C$ has degree at least
six. Thus $W$ has only word lengths one and two:
\[
                    W=s^{-1}w+W_2,
             \qquad w\in C_{15},\quad W_2\in T^2(s^{-1}\overline C).
             \tag{D.2.50g}
\]
The word-length-one part of \textup{(D.2.50f)} gives
$d_Cw=0$. Hence $w$ determines an actual integral homology
class $\omega=[w]\in H_{15}(L;\ZZ)$.

Define an integer-valued functional $\lambda$ on word-length
two and total cobar degree thirteen by
\[
 \lambda\big((s^{-1}c)(s^{-1}d)\big)=
 \begin{cases}
 A(c)B(d),& |c|=7,\ |d|=8,\\
 0,&\text{otherwise}.
 \end{cases}
 \tag{D.2.50h}
\]
Equations \textup{(D.2.50e)--(D.2.50f)} give
\[
                \lambda\big(2(uv-vu)\big)
                  =-2A(x)^2\equiv2\pmod4.
                  \tag{D.2.50i}
\]
The possible unshifted degree pairs in $W_2$ are exactly
$(7,9)$, $(8,8)$ and $(9,7)$. For the first pair,
$\lambda(d_1W_2)=0$ because $B(d_Cd)=(dB)(d)=0$.
For the middle pair its only possibly nonzero contribution
is $-A(d_Cc)B(d)=-4B(c)B(d)$, by \textup{(D.2.50d)}.
The last pair contributes zero for degree reasons. Therefore
\[
                            \lambda(d_1W_2)\equiv0\pmod4.
                            \tag{D.2.50j}
\]
There are no words of length three in the relevant degrees.
The length-two part of \textup{(D.2.50f)} is consequently
$d_2(s^{-1}w)+d_1W_2=2(uv-vu)$.
In the sign convention of \textup{(D.2.50c)}, desuspending
the first factor of degree seven gives
\[
 \lambda\big(d_2(s^{-1}w)\big)
             =-\langle A\smile B,w\rangle.
 \tag{D.2.50k}
\]
Combining \textup{(D.2.50i)--(D.2.50k)} proves
$\langle A\smile B,w\rangle\equiv2\pmod4$.
This is the first pairing in \textup{(D.2.50b)}. The sign
does not affect the value two modulo four.

\textbf{The half-product and its integral Bockstein.}
D.2.45 constructs $\theta$ with
\[
 a_7\smile\operatorname{red}_4b_8=i\theta,
 \qquad \beta_{\ZZ}\theta=b_8^2,
 \qquad i:\FF_2\longrightarrow\ZZ/4,\quad i(1)=2.
 \tag{D.2.50l}
\]
Its $\operatorname{Sq}^1H^{14}$ ambiguity evaluates to zero
on the reduction of every integral cycle. Thus the first
pairing implies the second pairing in \textup{(D.2.50b)}.
If $b_8^2=0$, the coefficient exact sequence would give an
integral lift of $\theta$. Every integral degree-fifteen
cohomology class on $L$ is torsion, since $L$ is rationally
acyclic and of finite type. Such a class evaluates to zero
on every integral cycle, contrary to
\textup{(D.2.50b)}. Hence $b_8^2\ne0$.
It has exact order two because it is the integral
Bockstein in \textup{(D.2.50l)}.

Finally, $H_{15}(L;\ZZ)$ is a finite torsion group.
Here is the needed mod-four lift of $\theta$ explicitly.
The half-product construction in D.2.45 supplies an integral
cochain $E\in C^{15}$ with $dE=2B^2$ and $E\bmod2$
representing $\theta$. The mod-two square vanishing proved
in D.2.42 allows an integral cochain $J\in C^{15}$ with
$dJ=B^2+2K$ for an integral cochain $K\in C^{16}$.
Consequently
\[
 \theta_4=[E-2J\bmod4]\in H^{15}(L;\ZZ/4),
 \qquad\operatorname{red}_2\theta_4=\theta,
 \quad d(E-2J)=-4K.
 \tag{D.2.50m}
\]
Thus $\theta$ evaluates to zero on every integral class
of order at most two: the evaluation of $\theta_4$ on such
a class lies in $\{0,2\}$ and reduces to zero modulo two.
The class $\omega$ from \textup{(D.2.50b)} therefore has
order at least four. The construction made $w$ an integral
cycle, with no claim that it is represented by a stable map
from a sphere.
\end{proof}
